\documentclass[11pt,letterpaper]{article}

\usepackage[margin=0.6in]{geometry}
\usepackage{enumitem}
\usepackage{amsmath}

\pagestyle{empty}
\setlength{\parindent}{0pt}

\begin{document}

\begin{center}
    {\Large \textbf{16.07 Dynamics --- Quiz 4: Kane Damper}}
\end{center}

\vspace{4pt}

Name: \rule{3.2in}{0.4pt}

\vspace{14pt}
The questions below refer to the rigid body with a Kane damper from Problem Set 4. Recall that the fluid applies a torque $C_d\left({}^{B}\boldsymbol{\omega}_d - {}^{B}\boldsymbol{\omega}\right)$ to the body, and that the total kinetic energy changes at the rate
\[
\dot T = -C_d\,\left\lVert {}^{B}\boldsymbol{\omega}_d - {}^{B}\boldsymbol{\omega} \right\rVert^2 .
\]

\vspace{10pt}
\begin{enumerate}[label=\textbf{\arabic*.}, leftmargin=*, itemsep=4pt, series=quiz]

\item Without a damper, spin about which principal axes is stable? Which of these does the Kane damper make unstable, and why?
\vspace{2in}

\item What happens to the rate of energy dissipation in the limit $C_d \to 0$? What about the limit $C_d \to \infty$?
\vspace{2in}

\item A body with principal inertia matrix $\mathbf{J} = \operatorname{diag}(2,\ 3,\ 4)$ contains a Kane damper with $J_d = 1$. It starts in an equilibrium spin at rate $\Omega$ about $\hat b_1$ (the minor axis), with the damper co-rotating, and is then given a tiny disturbance. About which axis, and at what angular velocity, does the body eventually spin?


\vspace*{\fill}

\end{enumerate}

\end{document}
